\chapter{Disease-Modifying Antirheumatic Drugs (DMARDs) for Arthritis}

\section*{Definition and Overview}
Disease-modifying antirheumatic drugs (DMARDs) are medicines that slow or stop the underlying disease process in inflammatory arthritis, rather than just relieving symptoms. They are the cornerstone of treatment for rheumatoid arthritis and other inflammatory joint conditions, such as psoriatic arthritis and ankylosing spondylitis. By reducing inflammation at its source, DMARDs prevent joint damage, preserve function, and improve long-term outcomes. There are two main groups: conventional synthetic DMARDs such as methotrexate, and biologic or targeted DMARDs such as tumour necrosis factor inhibitors. This chapter explains how DMARDs work, how they are used, and their risks.

\section*{Epidemiology}
Inflammatory arthritis affects hundreds of thousands of Australians, with rheumatoid arthritis affecting around 2\% of the population. Early treatment with DMARDs is now standard practice because it prevents permanent joint damage. Methotrexate is the most commonly used DMARD, and biologic medicines are used when conventional agents are not enough. The introduction of DMARDs, especially biologics, has dramatically improved outcomes, with fewer joint replacements and better function.

\section*{Aetiology and Pathophysiology}
\begin{itemize}
    \item \textbf{Inflammatory arthritis:} the immune system attacks the joints, causing inflammation, pain, swelling, and joint damage
    \item \textbf{Methotrexate:} reduces immune activity and inflammation; the most commonly used first-line DMARD
    \item \textbf{Other conventional DMARDs:} sulfasalazine, hydroxychloroquine, and leflunomide
    \item \textbf{Biologics:} antibodies that block specific inflammatory proteins, such as tumour necrosis factor (TNF), interleukins, and other targets
    \item \textbf{Targeted synthetic DMARDs:} oral medicines such as tofacitinib and baricitinib that block inflammatory signalling
    \item \textbf{Mechanism:} DMARDs interrupt the immune processes driving joint inflammation, preventing damage
\end{itemize}

\section*{Clinical Features}
\begin{itemize}
    \item DMARDs are started early after a diagnosis of inflammatory arthritis
    \item Methotrexate is taken weekly, usually by mouth or injection
    \item Effects develop over 6--12 weeks, so patients must continue treatment even before improvement
    \item Biologics are given by injection or infusion, often every 2--8 weeks
    \item DMARDs reduce pain, swelling, and stiffness and prevent joint damage
    \item Blood tests are needed regularly to monitor for side effects
\end{itemize}

\section*{Differential Diagnosis}
\begin{itemize}
    \item Conventional versus biologic DMARDs, chosen based on disease severity and response
    \item DMARDs versus simple pain relief and anti-inflammatory medicines, which treat symptoms only
    \item Different biologic targets, selected by the rheumatologist
    \item DMARD therapy versus other treatments for the specific type of arthritis
\end{itemize}

\section*{When to Seek Medical Care}
DMARDs are prescribed by a rheumatologist, usually after referral from a GP. Seek medical care for any new joint swelling, pain, or stiffness lasting more than a few weeks, and report any side effects or infections while taking DMARDs.

\textbf{Call 000 (or local emergency services) for:}
\begin{itemize}
    \item High fever, severe infection, or difficulty breathing while taking DMARDs
    \item Signs of a severe allergic reaction to a medicine
    \item Severe abdominal pain, unexplained bleeding, or sudden weakness
    \item Chest pain or breathlessness
\end{itemize}

\section*{Diagnosis and Investigations}
\begin{itemize}
    \item \textbf{Blood tests before starting:} full blood count, liver and kidney function, and hepatitis and tuberculosis screening for biologics
    \item \textbf{Regular monitoring:} blood counts and liver function during treatment
    \item \textbf{Assessment of disease activity:} joint examination and inflammatory markers
    \item \textbf{Imaging:} X-rays and ultrasound to assess joint damage
    \item \textbf{Review:} regular rheumatology review of response and side effects
\end{itemize}

\section*{Management}
\begin{itemize}
    \item \textbf{Methotrexate:} first-line treatment, with folic acid to reduce side effects
    \item \textbf{Combination therapy:} two or more conventional DMARDs when needed
    \item \textbf{Biologics:} added when conventional DMARDs fail to control the disease
    \item \textbf{Targeted synthetic DMARDs:} oral options for some people
    \item \textbf{Dose adjustment:} tailored by the rheumatologist to the response
    \item \textbf{Infection prevention:} vaccinations as recommended, and prompt treatment of infections
    \item \textbf{Long-term review:} ongoing monitoring of disease, function, and side effects
\end{itemize}

\section*{Complications}
\begin{itemize}
    \item Increased risk of infection from immune suppression
    \item Liver, kidney, or blood side effects, monitored by blood tests
    \item Methotrexate side effects including nausea and mouth ulcers
    \item Allergic or infusion reactions to biologics
    \item Uncontrolled disease with permanent joint damage if treatment is not effective
    \item Costs and the need for ongoing monitoring
\end{itemize}

\section*{Prognosis and Outcomes}
DMARDs have transformed the outlook for inflammatory arthritis. Early treatment controls inflammation, prevents joint damage, and maintains function, and many people achieve remission. Biologic medicines are highly effective for people who do not respond to conventional DMARDs. With monitoring and appropriate management of side effects, most people with inflammatory arthritis live active lives with far less disability than in the past.

\section*{Prevention and Self-Care}
\begin{itemize}
    \item Take DMARDs exactly as prescribed, including weekly methotrexate dosing
    \item Attend regular blood tests and rheumatology reviews
    \item Report infections, fevers, or side effects promptly
    \item Keep vaccinations up to date as advised
    \item Maintain joint-friendly exercise and a healthy lifestyle
    \item Use support services and patient education programs
    \item Discuss family planning with your doctor, as some DMARDs are not suitable in pregnancy
\end{itemize}

\section*{Sources and Further Reading}
The following reputable sources provide further information for healthcare students and patients seeking reliable, current advice about DMARDs for arthritis.

\begin{itemize}
    \item Arthritis Australia. \textit{DMARDs and arthritis treatment.} https://arthritisaustralia.com.au/
    \item Australian Rheumatology Association. \textit{Information about DMARDs.} https://rheumatology.org.au/
    \item Healthdirect Australia. \textit{Disease-modifying antirheumatic drugs.} https://www.healthdirect.gov.au/
    \item NHS. \textit{DMARDs: medicines for inflammatory arthritis.} https://www.nhs.uk/conditions/dmards/
    \item Mayo Clinic. \textit{Rheumatoid arthritis treatment: DMARDs.} https://www.mayoclinic.org/
    \item MSD Manual. \textit{Disease-modifying antirheumatic drugs.} https://www.msdmanuals.com/
\end{itemize}
